% lecture02.tex - 线性组合、张成与基
% 第 2 章: 线性组合、张成与基

\chapter{线性组合、张成与基}
\label{ch:span-basis}

\begin{abstract}
在本章中, 我们将探讨``线性组合 (linear combination)''、``张成 (span)''与``基 (basis)''这些核心概念。这些概念是深入理解线性代数的基础, 让我们明白向量 (vector) 是如何``填满''空间的。
\end{abstract}

% ============================================
\section{坐标的新视角}
% ============================================

在上一章中, 我们把向量坐标看作``移动指南''。但还有另一种非常重要的视角。

\begin{intuition}
当我们写 $\vv = \twovec{3}{-2}$ 时, 不要只想``向右走 3 个单位, 向下走 2 个单位'', 而要这样想:
\begin{itemize}
    \item 数字 3 是一个\textbf{标量 (scalar)}, 用来缩放向量 $\vi$
    \item 数字 $-2$ 是一个\textbf{标量}, 用来缩放向量 $\vj$
    \item 最终的向量是这两个缩放后向量的\textbf{和}
\end{itemize}
\end{intuition}

\begin{center}
\begin{tikzpicture}[scale=1.3]
    % 网格
    \draw[grid] (-1.5,-2.5) grid (4.5,1.5);

    % 坐标轴
    \draw[axis] (-1.5,0) -- (4.5,0) node[right] {$x$};
    \draw[axis] (0,-2.5) -- (0,1.5) node[above] {$y$};

    % 基向量 (单位向量)
    \draw[basis, -{Stealth[length=3mm]}] (0,0) -- (1,0) node[below right] {$\vi$};
    \draw[basis, -{Stealth[length=3mm]}] (0,0) -- (0,1) node[above left] {$\vj$};

    % 缩放后的向量
    \draw[vector, orange!70, very thick] (0,0) -- (3,0) node[midway, below] {$3\vi$};
    \draw[vector, green!60!black, very thick] (3,0) -- (3,-2) node[midway, right] {$-2\vj$};

    % 结果向量
    \draw[vector, blue!70, ultra thick] (0,0) -- (3,-2) node[below right] {$\vv = 3\vi - 2\vj$};
\end{tikzpicture}
\end{center}

% ============================================
\section{线性组合}
% ============================================

\begin{definition}[线性组合 (linear combination)]
两个向量 $\vv$ 与 $\vw$ 的\vocab{线性组合}为:
\[
a\vv + b\vw
\]
其中 $a$ 与 $b$ 为任意标量。

更一般地, $n$ 个向量 $\vv_1, \vv_2, \ldots, \vv_n$ 的线性组合为:
\[
c_1\vv_1 + c_2\vv_2 + \cdots + c_n\vv_n = \sum_{i=1}^{n} c_i\vv_i
\]
\end{definition}

\begin{intuition}
``线性 (linear)''这个名字从何而来? 如果固定其中一个标量, 让另一个标量变化, 所得向量的箭头终点会沿一条\textbf{直线}移动。

例如, 若固定 $b = 1$, 让 $a$ 变化, 则向量 $a\vv + \vw$ 会沿一条与 $\vv$ 平行的直线移动。
\end{intuition}

\begin{example}
设 $\vv = \twovec{1}{2}$, $\vw = \twovec{3}{-1}$。一些线性组合如下:
\begin{align*}
2\vv + 1\vw &= 2\twovec{1}{2} + 1\twovec{3}{-1} = \twovec{2}{4} + \twovec{3}{-1} = \twovec{5}{3} \\
-1\vv + 3\vw &= -\twovec{1}{2} + 3\twovec{3}{-1} = \twovec{-1}{-2} + \twovec{9}{-3} = \twovec{8}{-5} \\
0\vv + 0\vw &= \twovec{0}{0} = \vzero
\end{align*}
\end{example}

% ============================================
\section{张成}
% ============================================

\begin{definition}[张成 (span)]
一组向量的\vocab{张成}就是这些向量所有可能的线性组合构成的集合:
\[
\spn\{\vv_1, \vv_2, \ldots, \vv_n\} = \{c_1\vv_1 + c_2\vv_2 + \cdots + c_n\vv_n \mid c_i \in \R\}
\]
\end{definition}

\begin{intuition}
关键问题是: 给定一些向量, 只用两种基本运算 (加法与标量乘法), 我们能到达哪些向量?

答案就是: 这些向量的\textbf{张成}。
\end{intuition}

% -------------------------------------------
\subsection{二维空间中的张成}
% -------------------------------------------

\begin{theorem}
对于 $\Rtwo$ 中的两个向量 $\vv$ 与 $\vw$, 有以下三种可能情形:
\begin{enumerate}
    \item 若 $\vv$ 与 $\vw$ 不共线: $\spn\{\vv, \vw\} = \Rtwo$ (整个平面)
    \item 若 $\vv$ 与 $\vw$ 共线 (但非零): $\spn\{\vv, \vw\}$ 是一条直线
    \item 若两者均为零向量: $\spn\{\vv, \vw\} = \{\vzero\}$ (只有原点)
\end{enumerate}
\end{theorem}

\begin{center}
\begin{tikzpicture}[scale=0.9]
    % 情形 1: 整个平面
    \begin{scope}[shift={(0,0)}]
        \fill[blue!10] (-2,-2) rectangle (2,2);
        \draw[axis] (-2,0) -- (2,0);
        \draw[axis] (0,-2) -- (0,2);
        \draw[vector, red!70, very thick] (0,0) -- (1,0.5) node[right] {$\vv$};
        \draw[vector, green!60!black, very thick] (0,0) -- (-0.3,1) node[above] {$\vw$};
        \node at (0,-2.5) {情形 1: 整个平面};
    \end{scope}

    % 情形 2: 一条直线
    \begin{scope}[shift={(6,0)}]
        \draw[blue!30, very thick] (-2,-1) -- (2,1);
        \draw[axis] (-2,0) -- (2,0);
        \draw[axis] (0,-2) -- (0,2);
        \draw[vector, red!70, very thick] (0,0) -- (1,0.5) node[right] {$\vv$};
        \draw[vector, green!60!black, very thick] (0,0) -- (1.6,0.8) node[above] {$\vw$};
        \node at (0,-2.5) {情形 2: 一条直线};
    \end{scope}

    % 情形 3: 只有原点
    \begin{scope}[shift={(12,0)}]
        \draw[axis] (-2,0) -- (2,0);
        \draw[axis] (0,-2) -- (0,2);
        \fill[blue!70] (0,0) circle (3pt);
        \node at (0,-2.5) {情形 3: 只有原点};
    \end{scope}
\end{tikzpicture}
\end{center}

% -------------------------------------------
\subsection{三维空间中的张成}
% -------------------------------------------

\begin{theorem}
对于 $\Rthree$ 中的向量:
\begin{itemize}
    \item \textbf{一个非零向量:} 张成是一条\textbf{直线}
    \item \textbf{两个不共线的向量:} 张成是一个\textbf{平面}
    \item \textbf{三个不共面的向量:} 张成是整个 $\Rthree$
\end{itemize}
\end{theorem}

\begin{intuition}
想象三维空间中的两个向量。它们的线性组合构成一个过原点的平面。现在再加入一个不在该平面上的第三个向量, 就好比把这个平面在空间中``扫动''起来, 从而覆盖整个三维空间。
\end{intuition}

% ============================================
\section{线性无关与线性相关}
% ============================================

\begin{definition}[线性相关 (linear dependence)]
若一组向量中可以去掉其中一个而不改变张成, 则称这组向量\vocab{线性相关}, 换句话说, 其中至少有一个向量是``多余的''。

数学表述: 向量 $\vv_1, \ldots, \vv_n$ 线性相关, 当且仅当存在不全为零的标量 $c_1, \ldots, c_n$, 使得:
\[
c_1\vv_1 + c_2\vv_2 + \cdots + c_n\vv_n = \vzero
\]
\end{definition}

\begin{definition}[线性无关 (linear independence)]
若一组向量中没有一个是多余的---即每个向量都为张成增加了新的维度, 则称这组向量\vocab{线性无关}。

数学表述: 向量 $\vv_1, \ldots, \vv_n$ 线性无关, 当且仅当:
\[
c_1\vv_1 + c_2\vv_2 + \cdots + c_n\vv_n = \vzero \implies c_1 = c_2 = \cdots = c_n = 0
\]
\end{definition}

\begin{example}
考虑向量 $\vv = \twovec{2}{1}$ 与 $\vw = \twovec{4}{2}$。

这两个向量\textbf{线性相关}, 因为 $\vw = 2\vv$。我们可以写成:
\[
2\vv - \vw = \vzero \quad \text{即} \quad 2\twovec{2}{1} - \twovec{4}{2} = \twovec{0}{0}
\]
\end{example}

\begin{example}
向量 $\vv = \twovec{1}{0}$ 与 $\vw = \twovec{0}{1}$ 线性无关。

若 $a\vv + b\vw = \vzero$:
\[
a\twovec{1}{0} + b\twovec{0}{1} = \twovec{a}{b} = \twovec{0}{0}
\]
于是 $a = 0$ 且 $b = 0$。只有当所有系数都为零时, 才能得到零向量。
\end{example}

\begin{warning}
在 $n$ 维空间中, 至多有 $n$ 个向量能线性无关。如果向量个数超过 $n$, 它们一定线性相关。
\end{warning}

% ============================================
\section{基}
% ============================================

\begin{definition}[基 (basis)]
向量空间的一组\vocab{基}是满足以下两个条件的一组向量:
\begin{enumerate}
    \item \textbf{线性无关}
    \item \textbf{张成}整个空间
\end{enumerate}
\end{definition}

\begin{theorem}
$\Rtwo$ 的标准基为:
\[
\left\{ \vi = \twovec{1}{0}, \; \vj = \twovec{0}{1} \right\}
\]
$\Rthree$ 的标准基为:
\[
\left\{ \vi = \threevec{1}{0}{0}, \; \vj = \threevec{0}{1}{0}, \; \vk = \threevec{0}{0}{1} \right\}
\]
\end{theorem}

\begin{intuition}
基就像描述向量的一门``语言''。当我们说 $\twovec{3}{2}$ 时, 实际是在说``取 3 个第一个基向量, 再取 2 个第二个基向量''。

如果换一组基, 同一个向量的坐标也会不同---就像把一句话翻译成另一种语言。
\end{intuition}

% -------------------------------------------
\subsection{非标准基}
% -------------------------------------------

\begin{example}
以下集合也是 $\Rtwo$ 的一组基:
\[
\mathcal{B} = \left\{ \vb_1 = \twovec{1}{1}, \; \vb_2 = \twovec{1}{-1} \right\}
\]

\textbf{线性无关性的证明:} 若 $a\vb_1 + b\vb_2 = \vzero$:
\[
a\twovec{1}{1} + b\twovec{1}{-1} = \twovec{a+b}{a-b} = \twovec{0}{0}
\]
于是 $a + b = 0$ 且 $a - b = 0$, 解得 $a = b = 0$。

\textbf{张成整个平面:} 因为 $\Rtwo$ 中有两个线性无关的向量, 它们张成整个平面。
\end{example}

\begin{center}
\begin{tikzpicture}[scale=1.5]
    % 网格
    \draw[grid] (-2,-2) grid (2,2);

    % 坐标轴
    \draw[axis] (-2,0) -- (2,0) node[right] {$x$};
    \draw[axis] (0,-2) -- (0,2) node[above] {$y$};

    % 标准基
    \draw[basis, -{Stealth[length=3mm]}] (0,0) -- (1,0) node[below right] {$\vi$};
    \draw[basis, -{Stealth[length=3mm]}] (0,0) -- (0,1) node[above left] {$\vj$};

    % 非标准基
    \draw[vector, blue!70, very thick] (0,0) -- (1,1) node[above right] {$\vb_1$};
    \draw[vector, green!60!black, very thick] (0,0) -- (1,-1) node[below right] {$\vb_2$};
\end{tikzpicture}
\end{center}

% ============================================
\section{把向量看作点}
% ============================================

\begin{remark}
有时把向量看作一支箭头不如把它看作一个\textbf{点}---即向量箭头终点落到的那个点, 会更加方便。

\begin{itemize}
    \item 当思考单个具体的向量时: 把它想象成一支\textbf{箭头}
    \item 当思考一族向量时: 把它们想象成一些\textbf{点}
\end{itemize}
\end{remark}

\begin{center}
\begin{tikzpicture}[scale=1]
    % 左图: 作为箭头的向量
    \begin{scope}[shift={(0,0)}]
        \draw[axis] (-0.5,0) -- (3,0);
        \draw[axis] (0,-0.5) -- (0,3);
        \draw[vector, blue!70] (0,0) -- (1,2);
        \draw[vector, red!70] (0,0) -- (2,1);
        \draw[vector, green!60!black] (0,0) -- (2.5,2.5);
        \node at (1.5,-1) {作为箭头的向量};
    \end{scope}

    % 右图: 作为点的向量
    \begin{scope}[shift={(6,0)}]
        \draw[axis] (-0.5,0) -- (3,0);
        \draw[axis] (0,-0.5) -- (0,3);
        \fill[blue!70] (1,2) circle (3pt);
        \fill[red!70] (2,1) circle (3pt);
        \fill[green!60!black] (2.5,2.5) circle (3pt);
        \node at (1.5,-1) {作为点的向量};
    \end{scope}
\end{tikzpicture}
\end{center}

% ============================================
\section{核心概念总结}
% ============================================

\begin{summary}
\begin{description}
    \item[线性组合:] $c_1\vv_1 + c_2\vv_2 + \cdots + c_n\vv_n$, 其中系数为任意标量
    \item[张成:] 所有可能的线性组合构成的集合
    \item[线性无关:] 没有向量是多余的
    \item[线性相关:] 至少有一个向量可以去掉
    \item[基:] 线性无关且张成整个空间的向量组
\end{description}
\end{summary}

% ============================================
\section{练习}
% ============================================

\begin{exercise}
把向量 $\vv = \twovec{4}{6}$ 写成 $\vi$ 与 $\vj$ 的线性组合。
\end{exercise}

\begin{exercise}
以下向量是否线性无关? 请说明理由。
\[
\va = \twovec{1}{3}, \quad \vb = \twovec{2}{6}
\]
\end{exercise}

\begin{exercise}
以下向量是否线性无关?
\[
\vv_1 = \twovec{1}{0}, \quad \vv_2 = \twovec{1}{1}, \quad \vv_3 = \twovec{0}{1}
\]
\textit{提示: 在 $\Rtwo$ 中, 最多有几个线性无关的向量?}
\end{exercise}

\begin{exercise}
描述以下向量的张成:
\[
\vv = \threevec{1}{0}{0}, \quad \vw = \threevec{0}{1}{0}
\]
\end{exercise}

\begin{exercise}[挑战]
证明以下集合是 $\Rtwo$ 的一组基:
\[
\left\{ \twovec{2}{1}, \twovec{1}{3} \right\}
\]
然后求向量 $\twovec{5}{7}$ 在这组新基下的坐标。
\end{exercise}

\begin{exercise}
给定 $\Rthree$ 中的三个向量:
\[
\vv_1 = \threevec{1}{0}{1}, \quad \vv_2 = \threevec{0}{1}{1}, \quad \vv_3 = \threevec{1}{1}{2}
\]
这些向量是否线性无关? 它们的张成是什么?
\end{exercise}

\begin{problem}
证明: 若 $\{\vv_1, \vv_2\}$ 是 $\Rtwo$ 的一组基, 则 $\{2\vv_1, 3\vv_2\}$ 也是一组基。
\end{problem}

\begin{problem}
在 $\Rtwo$ 中能否找到三个线性无关的向量? 为什么?
\end{problem}
